\section{结论与后续工作}
\label{sec:conclusion}

本文完整介绍了 \sysname{}，一个面向 A 股上市公司的投研多智能体系统，覆盖财报分析、自动化估值建模与研究报告质量评估三个主选题，以及数据提取、产业链分析与研究报告纠错核查三个子选题。系统的出发点是一个朴素的观察：大语言模型在投研中最有价值的能力是读懂文档、组织论证与写作，最不可靠的能力是记住数字、计算数字与凭记忆补充事实。我们据此把系统拆成两半——模型负责读、抄、写、判断，程序负责校验、计算、拦截——并把这条分工落实为一组可由程序检查的不变量：每个数字带页码级来源并逐字回查；两条抽取通道冲突时不静默取舍；一张真值表驱动全部下游计算；文稿中的数字按展示精度在本次运行的产物或来源登记中找到出处；改写只能删减、不能新增或改动数字。

在此基础上，研报的生产被组织为五个上下文隔离、只靠文件交接的无头 Claude Code 会话与 11 道程序闸门；研报的评价被拆成程序判定的数据闭合与评审模型判定的分析质量，评分标准对生成方不可见，评审模型本身用重复评审、程序制造缺陷与人工复核三组证据校准。我们在真实年报上报告了抽取的覆盖与零错配结果，在流水线迭代中观察到研报质检得分的显著提升，并如实说明了这些证据的规模与局限。

我们认为，本工作中最具一般性的经验是：\textbf{与其要求模型更可靠，不如让模型的不可靠可以被发现。}系统中几乎每一次架构调整，都是把一个原本依赖模型自觉的性质，改成由程序在边界上强制。这一思路不局限于金融，适用于任何需要从长文档中提取事实、在事实之上写作、并对写作结果负责的场景。

\paragraph{后续工作。}按优先级，决赛前计划完成：
\begin{enumerate}
  \item \textbf{评测集扩充。}约 15 家公司的年报、季报与人工埋错的卖方研报，给出带置信区间的抽取与纠错指标。
  \item \textbf{对照消融。}在固定评分标准下，对资料先搜集、复核会话、深度专题三项设计做多次重复的对照实验；改进劣化算子使其跨全文同步删除内容与图表。
  \item \textbf{规则通道补全。}识别季报表头，使双通道交叉校验覆盖季报；补充股本变动表与政府补助口径。
  \item \textbf{跨家族评审。}引入不同家族的评审模型，检验同族偏好的影响。
  \item \textbf{执行器可替换性验证。}在国产模型执行器上重跑完整流水线，比较闸门通过率与质检得分。
  \item \textbf{交互前端。}网页前端与终端界面的运行中实时刷新，支持从研报中任一数字点击回溯到年报原文。
\end{enumerate}
